\documentclass[../main]{subfiles}
\begin{document}
\section{Dibujo Técnico Asistido}

\img{images/08/dibujo}{Dibujo en 3D.}
\info{Dibujo Asisitido}{
El dibujo asistido por computadora conocido como diseño asistido por ordenador, habitualmente citado como CAD por las siglas de su nombre en inglés computer-aided design es el uso de ordenadores para ayudar en la creación, modificación, análisis u optimización de un diseño.}

El software CAD se utiliza para aumentar la productividad del diseñador, mejorar la calidad del diseño, mejorar las comunicaciones a través de la documentación y crear una base de datos para la fabricación. La salida CAD a menudo se presenta en forma de archivos electrónicos para impresión, mecanizado u otras operaciones de fabricación. También se puede considerar al CAD como una técnica de dibujo.

Estas herramientas se pueden dividir básicamente en programas de dibujo 2D y de modelado 3D. Las herramientas de dibujo en 2D se basan en entidades geométricas vectoriales como puntos, líneas, arcos y polígonos, con las que se puede operar a través de una interfaz gráfica. Los modeladores en 3D añaden superficies y sólidos.
\newpage
\actividad{Leer, interpretar y responder}
\begin{enumerate}
    \item ¿Qué es CAD?
\end{enumerate}

\end{document}